\documentclass[12pt]{article}
\usepackage[utf8]{inputenc}
\usepackage[T1]{fontenc}
\usepackage[margin=1in]{geometry}
\usepackage{setspace}
\usepackage{xcolor}
\usepackage[colorlinks=true, urlcolor=blue!60!black, linkcolor=blue!60!black]{hyperref}
\setlength{\parindent}{0em}
\setlength{\parskip}{0.6em}
\renewcommand{\baselinestretch}{1.0}
\begin{document}
\begin{center}
{\LARGE \textit{United States v. Darby}, 312 U.S. 100 (1941)\par}
\vspace{0.5em}
{\large Case Brief --- Two-Page Edition\par}
\end{center}

\section{Memory Jogger}
Congress may bar interstate shipment of goods made under substandard labor conditions and regulate intrastate production for commerce; \textit{Hammer v.\ Dagenhart} is overruled and the Tenth Amendment is ``but a truism'' (\href{https://supreme.justia.com/cases/federal/us/312/100/}{312 U.S. 100}).

\section{Detailed Case Facts}
The Fair Labor Standards Act of 1938 set minimum wages and maximum hours for workers producing goods for interstate commerce, barred interstate shipment of goods made in violation, and required records. Georgia lumber manufacturer Fred Darby, who shipped much of his lumber out of state, was indicted for underpaying and overworking his mill workers, shipping the lumber, and failing to keep records.

\section{Procedural History}
Relying on \textit{Hammer}, the district court quashed the indictment; the United States appealed directly.

\section{Judicial Votes}
8--0. Stone wrote for a unanimous Court (one seat vacant); no separate opinions.

\section{Holding}
Congress may prohibit interstate shipment of goods made under substandard conditions, whatever its motive, and may directly regulate intrastate production for commerce as an appropriate means and because it substantially affects commerce. The Tenth Amendment adds no independent limit.

\section{Analysis of Opinions}
Under \textit{Gibbons}, the commerce power is complete and includes prohibition; motive is for Congress, and \textit{Hammer}, rejected by Holmes's ``powerful and now classic dissent,'' is overruled. Under \textit{McCulloch}, regulating production makes the shipment ban effective, and substandard wages fuel unfair interstate competition even from small producers. The Tenth Amendment is merely declaratory.

\section{Examples: Future Applications}
\textbf{Same-Side:} a shipment ban and direct penalties for electronics made under unsafe chemical exposure are valid. \textbf{Opposite-Side:} banning gun possession near schools is non-economic and invalid (\textit{Lopez}). \textbf{Fence-Sitter:} a purely local furniture shop that buys out-of-state supplies.

\section{Critique}
Barnett and Epstein say \textit{Darby} erased the line between commerce and production and gave Congress a police power; Ackerman and others defend it as recognizing a national economy and solving a race to the bottom.

\section{Key Quotations}
\begin{quote}
``The amendment states but a truism that all is retained which has not been surrendered. There is nothing in the history of its adoption to suggest that it was more than declaratory of the relationship between the national and state governments as it had been established by the Constitution before the amendment, or that its purpose was other than to allay fears that the new national government might seek to exercise powers not granted, and that the states might not be able to exercise fully their reserved powers.''
\end{quote}
\noindent --- Justice Stone, Opinion of the Court, 312 U.S. at 124.

\end{document}
